% ==============================================================================
% CHƯƠNG 4: TASK 2 — DẪN ĐƯỜNG AUV TRONG KHÔNG GIAN 3D
% ==============================================================================
\chapter{Task 2: Dẫn đường AUV trong không gian 3D}
\label{chap:task2}

\section{Mô hình hóa}
\label{sec:t2-mo-hinh}

Mỗi trạm khảo sát $i$ có tọa độ nguyên $(x_i, y_i, z_i)$ trong \ma{coords[100][3]}, với $z$ là
độ sâu. Cung $u \to v$ tồn tại khi $w(u, v) = $ \ma{weightMatrix[u][v]} $> 0$. Đi ngược áp suất
theo phương thẳng đứng tốn thêm năng lượng tỉ lệ với chênh lệch độ sâu:
\begin{equation}
  \NoiBat{c(u, v) = w(u, v)\left(1 + \frac{\abs{z_v - z_u}}{100}\right)}, \qquad
  g(n_{\text{mới}}) = g(n_{\text{cũ}}) + c(u, v).
  \label{eq:t2-cost}
\end{equation}
Vì $c(u, v) \ge w(u, v) > 0$, đồ thị có trọng số dương và A* luôn dừng. Heuristic là khoảng
cách 3D từ trạm hiện tại tới trạm đích theo ba chế độ của \eqref{eq:metric}:
\ma{mode} $= 1$ Manhattan, $2$ Euclid, $3$ Chebyshev; chế độ khác trả về \ma{nullptr}.

\begin{khungma}{Algo.cpp (trích findAUVPath)}{C++}
for (int u = 0; u < MAX_NODES; ++u) {
    for (int v = 0; v < MAX_NODES; ++v) {
        if (u == v || weightMatrix[u][v] <= 0) continue;
        double depthChange = std::fabs(static_cast<double>(coords[v][2] - coords[u][2]));
        graph.addEdge(u, v, weightMatrix[u][v] * (1.0 + depthChange / 100.0));
    }
    graph.h[u] = distance(metric, coords[u][0] - coords[goalPoint][0],
                          coords[u][1] - coords[goalPoint][1],
                          coords[u][2] - coords[goalPoint][2]);
}
\end{khungma}

\section{Phân tích lý thuyết ba heuristic}
\label{sec:t2-ly-thuyet}

Theo Định lý~\ref{dl:metric} với $d = 3$:
\[
  h_C \;\le\; h_E \;\le\; h_M \;\le\; 3\,h_C, \qquad h_M \le \sqrt{3}\, h_E.
\]
Tính chấp nhận được không phụ thuộc vào công thức metric mà vào \nhanmanh{quan hệ giữa đơn vị
năng lượng và đơn vị tọa độ}. Gọi
\begin{equation}
  \alpha = \min_{(u, v) \in E} \frac{c(u, v)}{d_E(u, v)}
  \label{eq:t2-alpha}
\end{equation}
là chi phí nhỏ nhất trên một đơn vị khoảng cách Euclid. Khi đó:
\begin{dsdung}
  \item Nếu $\alpha \ge 1$ thì $c(u, v) \ge d_E(u, v) \ge d_C(u, v)$, cộng bất đẳng thức tam
        giác ta được $h_E$ và $h_C$ đều \nhanmanh{nhất quán}; A* với hai chế độ này cho đường tối ưu.
  \item Trong cùng điều kiện, $h_E \ge h_C$ nên Euclid \emph{có nhiều thông tin hơn}: với heuristic
        nhất quán, tập đỉnh được mở rộng bởi $h_E$ là tập con của $h_C$ (trừ khác biệt do phá hòa)
        \cite{pearl1984}.
\end{dsdung}
\begin{dssai}
  \item $h_M$ có thể lớn hơn chi phí thật tới $\sqrt{3}$ lần nên chỉ chấp nhận được khi
        $\alpha \ge \sqrt{3}$ (hoặc khi mọi cung song song trục tọa độ). Ngược lại A* có thể trả
        đường không tối ưu, đổi lại thường mở rộng ít đỉnh nhất.
  \item Nếu $\alpha < 1$ thì cả ba metric đều có thể đánh giá quá; A* khi đó gần với tìm kiếm
        tham lam theo khoảng cách.
\end{dssai}

\section{Ví dụ tính tay: khi heuristic đánh giá quá}
\label{sec:t2-vi-du}

Mạng 5 trạm trong \ma{main.cpp} có tọa độ $0(0,0,0)$, $1(4,2,60)$, $2(5,0,10)$, $3(9,3,20)$,
$4(12,0,0)$ và các cạnh vô hướng $w$: $0\!-\!1 = 5$, $0\!-\!2 = 6$, $1\!-\!3 = 5$, $2\!-\!3 = 4$,
$3\!-\!4 = 4$, $2\!-\!4 = 9$. Theo \eqref{eq:t2-cost}: $c(0,1) = 8$, $c(0,2) = 6{,}6$,
$c(1,3) = 7$, $c(2,3) = 4{,}4$, $c(3,4) = 4{,}8$, $c(2,4) = 9{,}9$.

\begin{bangbieu}{Heuristic và chi phí thật còn lại tới trạm 4}{tab:t2-h}{C{1.2cm} C{2.6cm} C{1.8cm} C{1.8cm} C{1.8cm} Z}
  \hangtieude \thead{Trạm} & \thead{$(\Delta x, \Delta y, \Delta z)$} & \thead{$h_M$} & \thead{$h_E$} & \thead{$h_C$} & \thead{$h^\ast$} \\
  \midrule
  0 & $(12, 0, 0)$  & 12 & 12      & 12 & 15,8 \\
  1 & $(8, 2, 60)$  & \oxau 70 & \oxau 60,564 & \oxau 60 & 11,8 \\
  2 & $(7, 0, 10)$  & \oxau 17 & \oxau 12,207 & \oxau 10 & 9,2 \\
  3 & $(3, 3, 20)$  & \oxau 26 & \oxau 20,445 & \oxau 20 & 4,8 \\
  4 & $(0, 0, 0)$   & 0  & 0       & 0  & 0 \\
\end{bangbieu}

Ô đỏ là nơi heuristic lớn hơn $h^\ast$: ở mạng này $\alpha \approx 0{,}13 < 1$ (đạt tại cung $0 \to 1$; cung
$3 \to 4$ cũng dài $20{,}45$ đơn vị nhưng chỉ tốn $4{,}8$). Cả ba chế độ cho cùng một vết:

\begin{bangbieu}{Vết A* của Task 2 trên mạng mẫu (chế độ Euclid)}{tab:t2-vet}{C{0.9cm} C{4.4cm} Y}
  \hangtieude \thead{Vòng} & \thead{Mở rộng $(g; h; f)$} & \thead{Đưa vào \OPEN} \\
  \midrule
  1 & $0\ (0;\ 12;\ 12)$ & $1\ (8;\ 60{,}564;\ 68{,}564)$, $2\ (6{,}6;\ 12{,}207;\ 18{,}807)$ \\
  2 & $2\ (6{,}6;\ 12{,}207;\ 18{,}807)$ & $3\ (11;\ 20{,}445;\ 31{,}445)$, $4\ (16{,}5;\ 0;\ 16{,}5)$ \\
  3 & $4\ (16{,}5;\ 0;\ 16{,}5)$ & Đích được lấy ra, dừng \\
\end{bangbieu}

\begin{canhbao}[Đường tìm được không tối ưu]
  A* trả $0 \to 2 \to 4$ với $g = 16{,}5$, trong khi đường tối ưu là $0 \to 2 \to 3 \to 4$ với
  chi phí $6{,}6 + 4{,}4 + 4{,}8 = 15{,}8$. Trạm 3 có $f = 31{,}445 > 16{,}5$ vì $h_E(3) = 20{,}445$
  lớn gấp hơn 4 lần $h^\ast(3) = 4{,}8$, nên đích bị lấy ra trước khi trạm 3 kịp mở rộng. Đây đúng
  là hệ quả của Định lý~\ref{dl:toi-uu}: bảo đảm tối ưu chỉ có khi heuristic chấp nhận được.
\end{canhbao}

\section{So sánh thực nghiệm}
\label{sec:t2-thuc-nghiem}

\paragraph{Thiết lập} Mỗi bộ dữ liệu gồm $n$ trạm có tọa độ ngẫu nhiên trong
$[0, 100) \times [0, 100) \times [0, 60)$; mỗi trạm được nối với 4 trạm gần nhất (cạnh vô hướng).
Ở ba bộ \emph{liên kết} (C1--C3), $w(u, v) = \lceil d_E(u, v) \rceil$, nên $\alpha \ge 1$. Ở bộ
\emph{tách rời} (C4), $w$ là số nguyên ngẫu nhiên trong $[1, 9]$, không liên quan tới khoảng cách.
Đường đi luôn từ trạm $0$ tới trạm $n - 1$; chi phí tối ưu tính bằng Dijkstra. Số đỉnh mở rộng được
đọc từ biến đếm \ma{astar::expandedCount}; mọi lần chạy đều dưới 1~ms nên số đỉnh mở rộng là thước
đo hiệu năng đáng tin cậy hơn thời gian.

\begin{bangbieu}{So sánh ba heuristic 3D trên bốn bộ dữ liệu}{tab:t2-so-sanh}{C{1.2cm} C{1cm} C{1.8cm} Z Z Z Z}
  \hangtieude \thead{Bộ} & \thead{$n$} & \thead{Tối ưu} & \thead{Heuristic} & \thead{Chi phí A*} & \thead{Sai lệch} & \thead{Mở rộng} \\
  \midrule
  \multirow{3}{*}{C1} & \multirow{3}{*}{30} & \multirow{3}{*}{136,67}
     & Manhattan & \oxau 138,93 & \oxau +1,65\% & 10 \\
  & & & Euclid    & \otot 136,67 & 0\% & 11 \\
  & & & Chebyshev & \otot 136,67 & 0\% & 11 \\
  \midrule
  \multirow{3}{*}{C2} & \multirow{3}{*}{60} & \multirow{3}{*}{102,62}
     & Manhattan & \otot 102,62 & 0\% & 6 \\
  & & & Euclid    & \otot 102,62 & 0\% & 8 \\
  & & & Chebyshev & \otot 102,62 & 0\% & 10 \\
  \midrule
  \multirow{3}{*}{C3} & \multirow{3}{*}{90} & \multirow{3}{*}{129,52}
     & Manhattan & \oxau 129,80 & \oxau +0,22\% & 13 \\
  & & & Euclid    & \otot 129,52 & 0\% & 25 \\
  & & & Chebyshev & \otot 129,52 & 0\% & 35 \\
  \midrule
  \multirow{3}{*}{C4} & \multirow{3}{*}{60} & \multirow{3}{*}{26,73}
     & Manhattan & \oxau 39,78 & \oxau +48,82\% & 8 \\
  & & & Euclid    & \oxau 40,90 & \oxau +53,01\% & 10 \\
  & & & Chebyshev & \oxau 31,57 & \oxau +18,11\% & 8 \\
\end{bangbieu}
\nguonbang{Sinh bởi \ma{tools/report_data.cpp} (\ma{make report-data}), hạt giống 11--14.}

\begin{hangchiso}[3]
  \ChiSo[xanhla]{0\%}{sai lệch của Euclid và Chebyshev trên mọi bộ liên kết}
  \ChiSo[cam]{$-48\%$}{số đỉnh mở rộng của Manhattan so với Euclid ở C3}
  \ChiSo[do]{+53\%}{sai lệch lớn nhất khi chi phí tách rời tọa độ (C4)}
\end{hangchiso}

\section{Nhận xét}
\label{sec:t2-nhan-xet}

\SoSanh[xanhla][cam]{Euclid và Chebyshev}{
  \begin{dsdung}
    \item Luôn tối ưu khi $\alpha \ge 1$, đúng như phân tích ở \Muc{sec:t2-ly-thuyet}.
    \item Euclid sát $h^\ast$ hơn nên mở rộng ít hơn Chebyshev: 8 so với 10 (C2), 25 so với 35 (C3).
    \item Chebyshev yếu nhất do bỏ qua hai trong ba thành phần chênh lệch tọa độ.
  \end{dsdung}}
{Manhattan}{
  \begin{dsdung}
    \item Lớn nhất trong ba metric nên định hướng mạnh nhất, mở rộng ít đỉnh nhất (6, 10, 13).
  \end{dsdung}
  \begin{dssai}
    \item Không chấp nhận được khi $\alpha < \sqrt{3}$: lệch $+1{,}65\%$ (C1) và $+0{,}22\%$ (C3).
  \end{dssai}}

Khi chi phí không gắn với hình học (C4 và ví dụ ở \Muc{sec:t2-vi-du}), cả ba heuristic đều đánh
giá quá; Chebyshev, do nhỏ nhất, ít sai nhất ($+18\%$). Một cách khắc phục giữ nguyên dạng metric là
nhân heuristic với hệ số $\alpha$ ở \eqref{eq:t2-alpha}: $\alpha\, h_E$ và $\alpha\, h_C$ luôn nhất
quán. Chương trình không áp dụng hệ số này vì đề bài cố định công thức $h(n)$, nhưng đây là cách
nên dùng khi cần bảo đảm tối ưu với dữ liệu thực.
